\chapter{مبانی نظری و پیشینه پژوهش}
\label{chap:background}

\section{مدل ریزشبکه متصل به شبکه}
ریزشبکه مورد مطالعه همان تعریف فصل قبل است: بار محلی، منبع فتوولتائیک، باتری و یک اتصال به شبکه بالادست. در هر گام، توان بار و تولید \lr{PV} \textbf{برون‌زاد}اند؛ یعنی از بیرون مدل می‌آیند و کنترل‌کننده آن‌ها را عوض نمی‌کند. تصمیم فقط درباره استفاده از باتری و تبادل با شبکه است. این سطح، بهره‌برداری انرژی را توصیف می‌کند و جایگزین تحلیل گذرای ولتاژ، فرکانس یا حفاظت نیست.

نمادهای گام $t$ از این قرارند
\begin{itemize}
  \item $P_{\mathrm{L}}(t)$: توان بار بر حسب کیلووات
  \item $P_{\mathrm{PV}}(t)$: توان فتوولتائیکِ در دسترس، نامنفی
  \item $P_{\mathrm{G}}(t)$: توان خالص شبکه؛ مثبت برای خرید و منفی برای فروش
  \item $P_{\mathrm{c}}(t)$ و $P_{\mathrm{d}}(t)$: توان شارژ و دشارژ باتری، هر دو نامنفی
  \item $s(t)$: وضعیت شارژ، یعنی نسبت انرژی ذخیره‌شده به ظرفیت نامی
  \item $c_{\mathrm{b}}(t)$ و $c_{\mathrm{s}}(t)$: قیمت خرید و فروش برق
\end{itemize}

نمادگذاری توان و تفکیک لایه بهره‌برداری از لایه‌های اقتصادی گسترده‌تر با مسئله ریزشبکه متصل به شبکه هم‌راستا است \pcite{oladejo2019microgrid}. با این حال، در این گزارش فقط یک بهره‌بردار و یک اتصال شبکه وجود دارد.

\section{توازن توان و دینامیک باتری}
سمت چپ رابطه زیر منابع توان (خرید از شبکه، تولید فتوولتائیک و دشارژ) و سمت راست مصارف (بار و شارژ) است
\begin{equation}
  P_{\mathrm{G}}(t)+P_{\mathrm{PV}}(t)+P_{\mathrm{d}}(t)=P_{\mathrm{L}}(t)+P_{\mathrm{c}}(t)
  \label{eq:power-balance}
\end{equation}
اگر این برابری برقرار نباشد، \textbf{باقیمانده توازن} چنین تعریف می‌شود
\begin{equation}
  r_{\mathrm{G}}(t)=P_{\mathrm{G}}(t)+P_{\mathrm{PV}}(t)+P_{\mathrm{d}}(t)-P_{\mathrm{c}}(t)-P_{\mathrm{L}}(t)
  \label{eq:power-residual}
\end{equation}
مقدار صفر برای $r_{\mathrm{G}}$ یعنی توازن برقرار است. رفتار باتری با مدل انرژی گسسته تقریب زده می‌شود
\begin{equation}
  s(t+\Delta t)=s(t)+\frac{\eta_{\mathrm{c}}P_{\mathrm{c}}(t)\Delta t}{E_{\max}}-\frac{P_{\mathrm{d}}(t)\Delta t}{\eta_{\mathrm{d}}E_{\max}}
  \label{eq:soc-dynamics}
\end{equation}
که در آن $E_{\max}$ ظرفیت انرژی بر حسب کیلووات‌ساعت، $\Delta t=1\,\mathrm{h}$ طول گام، و $\eta_{\mathrm{c}}$ و $\eta_{\mathrm{d}}$ بازده شارژ و دشارژ هستند. حدهای ایمنی عبارت‌اند از
\begin{equation}
  s_{\min}\leq s(t)\leq s_{\max},\qquad
  0\leq P_{\mathrm{c}}(t),P_{\mathrm{d}}(t)\leq P_{\max}
  \label{eq:soc-bounds}
\end{equation}
برای یک باتری تک‌مبدله، شارژ و دشارژ هم‌زمان مجاز نیست. این \textbf{قید انحصار} به صورت زیر نوشته می‌شود
\begin{equation}
  P_{\mathrm{c}}(t)P_{\mathrm{d}}(t)=0
  \label{eq:exclusivity}
\end{equation}

\section{هزینه بهره‌برداری}
هزینه یک گام از سه بخش تشکیل می‌شود: هزینه خرید از شبکه، درآمد فروش به شبکه، و یک جریمه کوچک برای عبور انرژی از باتری
\begin{equation}
  g_{t}=c_{\mathrm{b}}(t)[P_{\mathrm{G}}(t)]_{+}\Delta t
  -c_{\mathrm{s}}(t)[-P_{\mathrm{G}}(t)]_{+}\Delta t
  +\kappa\bigl(P_{\mathrm{c}}(t)+P_{\mathrm{d}}(t)\bigr)\Delta t
  \label{eq:step-cost}
\end{equation}
عملگر $[z]_{+}=\max(0,z)$ بخش مثبت است؛ بنابراین فقط وقتی $P_{\mathrm{G}}$ مثبت است هزینه خرید و فقط وقتی منفی است درآمد فروش ظاهر می‌شود. وجود $\Delta t$ یکدستی واحدهای هزینه را حفظ می‌کند و رابطه را به گام یک‌ساعته محدود نمی‌کند. در مدل حاضر، $\kappa$ نماینده‌ای ساده‌شده از استهلاک عبوری باتری است، نه مدل الکتروشیمیایی واقعی. مدل‌های بهینه‌سازی بهره‌برداری نیز معمولاً هزینه انرژی را با حدهای توان و \lr{SOC} باتری همراه می‌کنند \pcite{elhendawi2018optimal}.

\section{یادگیری فیزیک‌آگاه و قید سخت}
در چارچوب \lr{PINN} متداول، شبکه خروجی آزاد می‌دهد و باقیمانده معادلات فیزیکی به تابع آموزش افزوده می‌شود \pcite{raissi2019pinns}. مزیت آن پیوند داده و مدل است. با این حال جریمه نرم، چون ضریب محدودی دارد، به‌تنهایی تضمین نمی‌کند که خروجی نهایی دقیقاً قید را رعایت کند. در بهره‌برداری انرژی حتی نقض کوچک \lr{SOC} یا توان با علامت نامعتبر می‌تواند فرمان را غیرقابل اجرا کند.

	extbf{قید سخت} یعنی قید در خود معماری اعمال شود، نه فقط در زیان. پژوهش‌ها نشان داده‌اند لایه‌های تحلیلی یا نرافکنی می‌توانند بعضی قیود را دقیقاً تضمین کنند \pcite{lu2021hardconstraints,chen2024hardlinear}. در این کار همین ایده برای قیود جبری ریزشبکه به کار می‌رود. روش حاضر یک \lr{PINN} حل‌کننده معادله دیفرانسیل جزئی نیست؛ تعبیر دقیق‌تر آن یادگیری فیزیک‌تعبیه‌شده است: فیزیک در نگاشت خروجی نشسته است، نه فقط در جریمه.

\section{جایگاه کار در پیشینه}
\subsection{مدیریت انرژی مبتنی بر بهینه‌سازی}
روش‌های کلاسیک، بهره‌برداری را به‌صورت برنامه‌ریزی خطی یا آمیخته عدد صحیح، برنامه‌ریزی پویا، یا کنترل پیش‌بین مدل صورت‌بندی می‌کنند. 	extbf{برنامه‌ریزی پویا} تصمیم جاری را با رابطه بلمن به هزینه آینده پیوند می‌دهد \pcite{bertsekas2017dp}. 	extbf{برنامه‌ریزی خطی آمیخته عدد صحیح} (\lr{MILP}) برای تصمیم‌های گسسته مناسب است و 	extbf{کنترل پیش‌بین مدل} (\lr{MPC}) افق لغزان را بازبهینه می‌کند. مطالعات موردی ریزشبکه نشان داده‌اند که \lr{MPC/MILP} می‌تواند هزینه و قیود بهره‌برداری را در یک صورت‌بندی صریح ترکیب کند \pcite{parisio2014mpc,parisio2014experimental,elhendawi2018optimal}. این روش‌ها مرجع‌هایی تفسیرپذیر فراهم می‌کنند، اما هزینه حل آن‌ها به طول افق، تفکیک شبکه حالت، سناریوها و جزئیات مدل وابسته است. برنامه‌ریزی پویای این گزارش نیز مرجع گسسته است و ادعای بهینگی پیوسته ندارد.

\subsection{مدیریت انرژی مبتنی بر یادگیری}
یادگیری تقویتی، سیاست کنترل را از تعامل با محیط می‌آموزد و شبکه‌های عصبی نیز می‌توانند به‌صورت جانشین نظارت‌شده، نگاشت خروجی یک حل‌کننده را تقریب بزنند. مرور کاربرد یادگیری تقویتی در پاسخ‌گویی بار، ظرفیت این روش‌ها و در عین حال دشواری ارزیابی و تعمیم را نشان می‌دهد \pcite{vazquez2019rl}. در این پروژه مسیر دوم دنبال می‌شود: برچسب‌ها از برنامه‌ریزی پویا می‌آیند و شبکه با بهینه‌ساز آدام آموزش می‌بیند \pcite{kingma2015adam}. مسئله مشترک آن است که خروجی آزاد یا سیاست یادگرفته‌شده، بدون سازوکار اضافی لزوماً توازن توان، حدود \lr{SOC} و انحصار متقابل شارژ و دشارژ را رعایت نمی‌کند؛ بنابراین خطای تقلید کم به‌تنهایی گواه امکان‌پذیری نیست.

\subsection{اعمال قید در کنترل‌کننده‌های یادگرفته‌شده}
در یادگیری تقویتی ایمن، مسئله را می‌توان به‌شکل فرایند تصمیم‌گیری مارکوف مقید نوشت \pcite{altman1999cmdp} و قید را با بهینه‌سازی سیاست مقید یا توابع لیاپانوف وارد کرد \pcite{achiam2017cpo,chow2018lyapunov}. بنچمارک‌های اکتشاف ایمن نیز تأکید می‌کنند که پاداش و ایمنی باید جداگانه گزارش شوند \pcite{ray2019safexp}. روش‌های جریمه‌ای ساده‌ترند، اما رعایت قید به وزن جریمه و کیفیت آموزش وابسته می‌ماند \pcite{garcia2015survey}.

راه دیگر، محدود کردن مستقیم فضای خروجی است. \lr{OptNet} یک برنامه درجه‌دو تفکیک‌پذیر را به‌صورت لایه شبکه وارد می‌کند \pcite{amos2017optnet} و شبکه‌های محدب نسبت به ورودی، ساختاری را فراهم می‌کنند که بهینه‌سازی روی برخی ورودی‌ها محدب بماند \pcite{amos2017icnn}. معماری‌های قید سخت نیز قید را با بازنپارامتری یا لایه تحلیلی اعمال می‌کنند \pcite{lu2021hardconstraints,chen2024hardlinear}. انتخاب حاضر از این خانواده است، اما حل‌کننده بهینه‌سازی تفکیک‌پذیر یا \lr{ICNN} نیست: یک نگاشت بسته و بدون پارامتر، تنها فرمان باتری را به چهار کمیت قابل اجرا تبدیل می‌کند. این سادگی مخصوص قیود جبری مدل حاضر است و برتری عمومی نسبت به روش‌های ایمن دیگر ادعا نمی‌شود.

\subsection{یادگیری فیزیک‌آگاه در سامانه‌های انرژی}
در شبکه‌های \lr{PINN} متعارف، باقیمانده معادله دیفرانسیل در تابع آموزش ظاهر می‌شود \pcite{raissi2019pinns} و چارچوب گسترده‌تر یادگیری فیزیک‌آگاه، ورود دانش فیزیکی از راه داده، زیان یا معماری را در بر می‌گیرد \pcite{karniadakis2021piml}. با این حال، مسئله این گزارش حل یک \lr{PDE} نیست. توازن توان و گذار ساعتی باتری روابط جبری و گسسته‌اند و در لایه خروجی دقیقاً محاسبه می‌شوند، نه آن‌که باقیمانده آن‌ها فقط کمینه شود. ازاین‌رو، اصطلاح دقیق‌تر برای روش حاضر «شبکه فیزیک‌تعبیه‌شده» است؛ صفر شدن نقض‌های تعریف‌شده خاصیت طراحی این نگاشت است، نه کشف آماری. جدول~\ref{tab:approach-comparison} نقش هر رویکرد را در این گزارش، نه لزوماً در یک آزمایش چهارطرفه، خلاصه می‌کند.

\begin{table}[H]
  \centering
  \caption{مقایسه رویکردهای تولید برنامه انرژی}
  \label{tab:approach-comparison}
  \begin{tabular}{>{\centering\arraybackslash}p{0.22\textwidth}>{\centering\arraybackslash}p{0.22\textwidth}>{\centering\arraybackslash}p{0.22\textwidth}>{\centering\arraybackslash}p{0.22\textwidth}}
    \toprule
    \textbf{رویکرد} & \textbf{داده یا مدل لازم} & \textbf{وضعیت قیدها} & \textbf{نقش در این پروژه} \\
    \midrule
    برنامه‌ریزی پویا & مدل و افق برون‌زاد & مقید، تا دقت شبکه حالت & مرجع اقتصادی \\
    شبکه داده‌محور آزاد & نمونه برچسب‌دار & تضمینی ندارد & \lr{baseline} یادگیری \\
    \lr{PINN} با جریمه نرم & نمونه و باقیمانده قانون & وابسته به وزن جریمه & فقط انگیزه روش‌شناختی؛ در این گزارش آموزش داده نشد \\
    شبکه فیزیک‌تعبیه‌شده & نمونه و نگاشت فیزیکی & سخت، برای قیود مدل‌شده & مدل پیشنهادی \\
    \bottomrule
  \end{tabular}
\end{table}

\section{مرجع برنامه‌ریزی پویا}
اگر شبکه فقط هزینه همان ساعت را کمینه کند، یک قاعده 	extbf{حریصانه} می‌آموزد و آینده باتری را نمی‌بیند. برای جلوگیری از این کار، برنامه مرجع با برنامه‌ریزی پویا ساخته می‌شود. حالت برابر \lr{SOC} و فرمان برابر توان علامت‌دار باتری است. اگر $V_{t}(s)$ کمینه هزینه از زمان $t$ به بعد در وضعیت $s$ باشد، رابطه بلمن چنین است
\begin{equation}
  V_{t}(s)=\min_{u\in\mathcal{U}(s)}\bigl\{g_{t}(s,u)+V_{t+1}\bigl(f(s,u)\bigr)\bigr\}
  \label{eq:bellman}
\end{equation}
که $f$ همان دینامیک~\eqref{eq:soc-dynamics} و $\mathcal{U}(s)$ مجموعه فرمان‌های مجاز در \lr{SOC} فعلی است. در پیاده‌سازی حاضر شرط نهایی $V_{T}(s)=0$ برای همه نقاط شبکه وضعیت در نظر گرفته شده است؛ بنابراین ارزش بازیافت یا هدف نهایی جداگانه‌ای برای باتری اعمال نمی‌شود. \lr{SOC} در مرجع روی شبکه‌ای با گام $0{.}01$ گسسته می‌شود؛ بنابراین مرجع، تقریب عددیِ یک حل‌کننده افق کامل با آینده‌نگری کامل است، نه بهینه دقیق پیوسته.

\section{جمع‌بندی فصل}
این فصل قیود اصلی و شاخص هزینه را مشخص کرد. تفاوت اساسی شبکه فیزیک‌تعبیه‌شده با مدل آزاد آن است که امکان‌پذیریِ قیود مدل‌شده تنها یک معیار پس از آموزش نیست؛ بخشی از تعریف خروجی است. فصل بعد نحوه ساخت داده، معماری و آزمون بازگشتی را دقیق می‌کند.
